% Studies-in-Algebra-and-Group-Theory - Lecture 16: Real and Quaternionic Representations
% Author: S. D. V. Stephens
% Date: 2026-03-10

\documentclass{report}
\input{~/university/preamble.tex}
% preamble-darkmode.tex
% Dark mode extension for lecture notes
% Usage: % preamble-darkmode.tex
% Dark mode extension for lecture notes
% Usage: % preamble-darkmode.tex
% Dark mode extension for lecture notes
% Usage: \input{~/university/preamble-darkmode.tex} AFTER preamble.tex
%
% This provides:
% - Dark background with light text
% - Material Design color palette
% - Ocean blue gradient colors (p1-p9)
% - Enhanced TikZ/PGFPlots for dark backgrounds
% - qq tcolorbox environment
% - tkz-euclide for geometric constructions

% ===========================================
% DARK MODE PACKAGE
% ===========================================
\usepackage{darkmode}
\enabledarkmode

% ===========================================
% ADDITIONAL PACKAGES
% ===========================================
\usepackage{changepage}
\usepackage{ulem}
\usepackage{colortbl}
\usepackage{tkz-euclide}
\usepackage{capt-of}

% ===========================================
% EXTENDED TIKZ LIBRARIES
% ===========================================
\usetikzlibrary{patterns,3d,backgrounds}
\usepgfplotslibrary{groupplots}

% Upgrade pgfplots compatibility
\pgfplotsset{compat=1.18}

% ===========================================
% OCEAN BLUE GRADIENT (p1-p9)
% ===========================================
\definecolor{p1}{HTML}{caf0f8}
\definecolor{p2}{HTML}{ade8f4}
\definecolor{p3}{HTML}{90e0ef}
\definecolor{p4}{HTML}{48cae4}
\definecolor{p5}{HTML}{00b4d8}
\definecolor{p6}{HTML}{0096c7}
\definecolor{p7}{HTML}{0077b6}
\definecolor{p8}{HTML}{023e8a}
\definecolor{p9}{HTML}{03045e}

% Dark background color
\definecolor{pag}{HTML}{293133}

% ===========================================
% MATERIAL DESIGN COLORS
% ===========================================
% Note: myred, myyellow already defined in main preamble
% These are additional/alternate versions
\definecolor{matred}{HTML}{f44336}
\definecolor{matpink}{HTML}{e81e63}
\definecolor{matpurple}{HTML}{9c27b0}
\definecolor{matdeeppurple}{HTML}{673ab7}
\definecolor{matindigo}{HTML}{3f51b5}
\definecolor{matblue}{HTML}{2196f3}
\definecolor{matlightblue}{HTML}{03a9f4}
\definecolor{matcyan}{HTML}{00bcd4}
\definecolor{matteal}{HTML}{009688}
\definecolor{matgreen}{HTML}{4caf50}
\definecolor{matlightgreen}{HTML}{8bc34a}
\definecolor{matlime}{HTML}{cddc39}
\definecolor{matyellow}{HTML}{ffeb3b}
\definecolor{matamber}{HTML}{ffc107}
\definecolor{matorange}{HTML}{ff9800}
\definecolor{matdeeporange}{HTML}{ff5722}
\definecolor{matbrown}{HTML}{795548}
\definecolor{matgray}{HTML}{9e9e9e}
\definecolor{matbluegray}{HTML}{607d8b}

% ===========================================
% COLORLET FOR TIKZ
% ===========================================
\colorlet{xcol}{blue!60!black}

% ===========================================
% DARK MODE TCOLORBOX
% ===========================================
\newtcolorbox{qq}[2][]{%
    boxrule=0.75pt,
    sharp corners,
    colframe=white,
    colback=pag,
    coltext=white,
    #1,
}

% Dark definition box
\newtcolorbox{darkdfn}[2][]{%
    enhanced,
    boxrule=1pt,
    sharp corners,
    colframe=p5,
    colback=pag,
    coltext=white,
    coltitle=white,
    fonttitle=\bfseries,
    title=#2,
    #1,
}

% Dark theorem box
\newtcolorbox{darkthm}[2][]{%
    enhanced,
    boxrule=1pt,
    sharp corners,
    colframe=matpurple,
    colback=pag,
    coltext=white,
    coltitle=white,
    fonttitle=\bfseries,
    title=#2,
    #1,
}

% ===========================================
% TIKZ 3D FIX
% ===========================================
% Small fix for canvas is xy plane at z
% https://tex.stackexchange.com/a/48776/121799
\makeatletter
\tikzoption{canvas is xy plane at z}[]{%
    \def\tikz@plane@origin{\pgfpointxyz{0}{0}{#1}}%
    \def\tikz@plane@x{\pgfpointxyz{1}{0}{#1}}%
    \def\tikz@plane@y{\pgfpointxyz{0}{1}{#1}}%
    \tikz@canvas@is@plane}
\makeatother

% ===========================================
% CONDITIONAL LINE TIKZSET
% ===========================================
\tikzset{
    conditional line/.style 2 args={
        /utils/exec={\pgfmathsetmacro{\xcoordA}{#1}},
        /utils/exec={\pgfmathsetmacro{\xcoordB}{#2}},
        insert path={
            \ifdim\xcoordA pt>\xcoordB pt
            (\xcoordA,0) -- (\xcoordB,0)
            \fi
        },
    },
}

% ===========================================
% PGFPLOTS DARK MODE DEFAULTS
% ===========================================
\pgfplotsset{
    dark axis/.style={
        axis line style={white},
        tick style={white},
        label style={white},
        grid style={pag!70},
        every axis label/.append style={white},
        every tick label/.append style={white},
    },
}

% ===========================================
% TIKZ EXTERNALIZATION (for fast compilation)
% ===========================================
% This creates .md5 cache files for figures
% Requires: pdflatex -shell-escape
%
% To enable, add AFTER loading this file:
%   \enabletikzexternalize
%
\newcommand{\enabletikzexternalize}{%
  \usetikzlibrary{external}%
  \tikzexternalize[prefix=figures/]%
}

% ===========================================
% CONVENIENT SHORTCUTS FOR DARK PLOTS
% ===========================================
\newcommand{\darkplot}[1]{%
    \begin{tikzpicture}
    \begin{axis}[
        dark axis,
        grid=both,
        major grid style={line width=.2pt,draw=pag!70},
        #1
    ]
}


% ===========================================
% QUICK GEOMETRY MACROS (tkz-euclide)
% ===========================================
% Draw a point: \point{name}{x}{y}
\newcommand{\gpoint}[3]{\tkzDefPoint(#2,#3){#1}}

% Draw line through points: \gline{A}{B}
\newcommand{\gline}[2]{\tkzDrawLine[color=p4](#1,#2)}

% Draw circle: \gcircle{center}{radius}
\newcommand{\gcircle}[2]{\tkzDrawCircle[color=p5](#1,#2)}

% Mark angle: \gangle{A}{B}{C}
\newcommand{\gangle}[3]{\tkzMarkAngle[color=p3,size=0.5](#1,#2,#3)}

% ===========================================
% USAGE EXAMPLE
% ===========================================
% In your document:
%
% \begin{tikzpicture}
%   \begin{axis}[dark axis, ...]
%     \addplot[p4, thick] {sin(deg(x))};
%   \end{axis}
% \end{tikzpicture}
%
% Or with qq box:
% \begin{qq}{My Title}
%   Content here in dark box
% \end{qq}
 AFTER preamble.tex
%
% This provides:
% - Dark background with light text
% - Material Design color palette
% - Ocean blue gradient colors (p1-p9)
% - Enhanced TikZ/PGFPlots for dark backgrounds
% - qq tcolorbox environment
% - tkz-euclide for geometric constructions

% ===========================================
% DARK MODE PACKAGE
% ===========================================
\usepackage{darkmode}
\enabledarkmode

% ===========================================
% ADDITIONAL PACKAGES
% ===========================================
\usepackage{changepage}
\usepackage{ulem}
\usepackage{colortbl}
\usepackage{tkz-euclide}
\usepackage{capt-of}

% ===========================================
% EXTENDED TIKZ LIBRARIES
% ===========================================
\usetikzlibrary{patterns,3d,backgrounds}
\usepgfplotslibrary{groupplots}

% Upgrade pgfplots compatibility
\pgfplotsset{compat=1.18}

% ===========================================
% OCEAN BLUE GRADIENT (p1-p9)
% ===========================================
\definecolor{p1}{HTML}{caf0f8}
\definecolor{p2}{HTML}{ade8f4}
\definecolor{p3}{HTML}{90e0ef}
\definecolor{p4}{HTML}{48cae4}
\definecolor{p5}{HTML}{00b4d8}
\definecolor{p6}{HTML}{0096c7}
\definecolor{p7}{HTML}{0077b6}
\definecolor{p8}{HTML}{023e8a}
\definecolor{p9}{HTML}{03045e}

% Dark background color
\definecolor{pag}{HTML}{293133}

% ===========================================
% MATERIAL DESIGN COLORS
% ===========================================
% Note: myred, myyellow already defined in main preamble
% These are additional/alternate versions
\definecolor{matred}{HTML}{f44336}
\definecolor{matpink}{HTML}{e81e63}
\definecolor{matpurple}{HTML}{9c27b0}
\definecolor{matdeeppurple}{HTML}{673ab7}
\definecolor{matindigo}{HTML}{3f51b5}
\definecolor{matblue}{HTML}{2196f3}
\definecolor{matlightblue}{HTML}{03a9f4}
\definecolor{matcyan}{HTML}{00bcd4}
\definecolor{matteal}{HTML}{009688}
\definecolor{matgreen}{HTML}{4caf50}
\definecolor{matlightgreen}{HTML}{8bc34a}
\definecolor{matlime}{HTML}{cddc39}
\definecolor{matyellow}{HTML}{ffeb3b}
\definecolor{matamber}{HTML}{ffc107}
\definecolor{matorange}{HTML}{ff9800}
\definecolor{matdeeporange}{HTML}{ff5722}
\definecolor{matbrown}{HTML}{795548}
\definecolor{matgray}{HTML}{9e9e9e}
\definecolor{matbluegray}{HTML}{607d8b}

% ===========================================
% COLORLET FOR TIKZ
% ===========================================
\colorlet{xcol}{blue!60!black}

% ===========================================
% DARK MODE TCOLORBOX
% ===========================================
\newtcolorbox{qq}[2][]{%
    boxrule=0.75pt,
    sharp corners,
    colframe=white,
    colback=pag,
    coltext=white,
    #1,
}

% Dark definition box
\newtcolorbox{darkdfn}[2][]{%
    enhanced,
    boxrule=1pt,
    sharp corners,
    colframe=p5,
    colback=pag,
    coltext=white,
    coltitle=white,
    fonttitle=\bfseries,
    title=#2,
    #1,
}

% Dark theorem box
\newtcolorbox{darkthm}[2][]{%
    enhanced,
    boxrule=1pt,
    sharp corners,
    colframe=matpurple,
    colback=pag,
    coltext=white,
    coltitle=white,
    fonttitle=\bfseries,
    title=#2,
    #1,
}

% ===========================================
% TIKZ 3D FIX
% ===========================================
% Small fix for canvas is xy plane at z
% https://tex.stackexchange.com/a/48776/121799
\makeatletter
\tikzoption{canvas is xy plane at z}[]{%
    \def\tikz@plane@origin{\pgfpointxyz{0}{0}{#1}}%
    \def\tikz@plane@x{\pgfpointxyz{1}{0}{#1}}%
    \def\tikz@plane@y{\pgfpointxyz{0}{1}{#1}}%
    \tikz@canvas@is@plane}
\makeatother

% ===========================================
% CONDITIONAL LINE TIKZSET
% ===========================================
\tikzset{
    conditional line/.style 2 args={
        /utils/exec={\pgfmathsetmacro{\xcoordA}{#1}},
        /utils/exec={\pgfmathsetmacro{\xcoordB}{#2}},
        insert path={
            \ifdim\xcoordA pt>\xcoordB pt
            (\xcoordA,0) -- (\xcoordB,0)
            \fi
        },
    },
}

% ===========================================
% PGFPLOTS DARK MODE DEFAULTS
% ===========================================
\pgfplotsset{
    dark axis/.style={
        axis line style={white},
        tick style={white},
        label style={white},
        grid style={pag!70},
        every axis label/.append style={white},
        every tick label/.append style={white},
    },
}

% ===========================================
% TIKZ EXTERNALIZATION (for fast compilation)
% ===========================================
% This creates .md5 cache files for figures
% Requires: pdflatex -shell-escape
%
% To enable, add AFTER loading this file:
%   \enabletikzexternalize
%
\newcommand{\enabletikzexternalize}{%
  \usetikzlibrary{external}%
  \tikzexternalize[prefix=figures/]%
}

% ===========================================
% CONVENIENT SHORTCUTS FOR DARK PLOTS
% ===========================================
\newcommand{\darkplot}[1]{%
    \begin{tikzpicture}
    \begin{axis}[
        dark axis,
        grid=both,
        major grid style={line width=.2pt,draw=pag!70},
        #1
    ]
}


% ===========================================
% QUICK GEOMETRY MACROS (tkz-euclide)
% ===========================================
% Draw a point: \point{name}{x}{y}
\newcommand{\gpoint}[3]{\tkzDefPoint(#2,#3){#1}}

% Draw line through points: \gline{A}{B}
\newcommand{\gline}[2]{\tkzDrawLine[color=p4](#1,#2)}

% Draw circle: \gcircle{center}{radius}
\newcommand{\gcircle}[2]{\tkzDrawCircle[color=p5](#1,#2)}

% Mark angle: \gangle{A}{B}{C}
\newcommand{\gangle}[3]{\tkzMarkAngle[color=p3,size=0.5](#1,#2,#3)}

% ===========================================
% USAGE EXAMPLE
% ===========================================
% In your document:
%
% \begin{tikzpicture}
%   \begin{axis}[dark axis, ...]
%     \addplot[p4, thick] {sin(deg(x))};
%   \end{axis}
% \end{tikzpicture}
%
% Or with qq box:
% \begin{qq}{My Title}
%   Content here in dark box
% \end{qq}
 AFTER preamble.tex
%
% This provides:
% - Dark background with light text
% - Material Design color palette
% - Ocean blue gradient colors (p1-p9)
% - Enhanced TikZ/PGFPlots for dark backgrounds
% - qq tcolorbox environment
% - tkz-euclide for geometric constructions

% ===========================================
% DARK MODE PACKAGE
% ===========================================
\usepackage{darkmode}
\enabledarkmode

% ===========================================
% ADDITIONAL PACKAGES
% ===========================================
\usepackage{changepage}
\usepackage{ulem}
\usepackage{colortbl}
\usepackage{tkz-euclide}
\usepackage{capt-of}

% ===========================================
% EXTENDED TIKZ LIBRARIES
% ===========================================
\usetikzlibrary{patterns,3d,backgrounds}
\usepgfplotslibrary{groupplots}

% Upgrade pgfplots compatibility
\pgfplotsset{compat=1.18}

% ===========================================
% OCEAN BLUE GRADIENT (p1-p9)
% ===========================================
\definecolor{p1}{HTML}{caf0f8}
\definecolor{p2}{HTML}{ade8f4}
\definecolor{p3}{HTML}{90e0ef}
\definecolor{p4}{HTML}{48cae4}
\definecolor{p5}{HTML}{00b4d8}
\definecolor{p6}{HTML}{0096c7}
\definecolor{p7}{HTML}{0077b6}
\definecolor{p8}{HTML}{023e8a}
\definecolor{p9}{HTML}{03045e}

% Dark background color
\definecolor{pag}{HTML}{293133}

% ===========================================
% MATERIAL DESIGN COLORS
% ===========================================
% Note: myred, myyellow already defined in main preamble
% These are additional/alternate versions
\definecolor{matred}{HTML}{f44336}
\definecolor{matpink}{HTML}{e81e63}
\definecolor{matpurple}{HTML}{9c27b0}
\definecolor{matdeeppurple}{HTML}{673ab7}
\definecolor{matindigo}{HTML}{3f51b5}
\definecolor{matblue}{HTML}{2196f3}
\definecolor{matlightblue}{HTML}{03a9f4}
\definecolor{matcyan}{HTML}{00bcd4}
\definecolor{matteal}{HTML}{009688}
\definecolor{matgreen}{HTML}{4caf50}
\definecolor{matlightgreen}{HTML}{8bc34a}
\definecolor{matlime}{HTML}{cddc39}
\definecolor{matyellow}{HTML}{ffeb3b}
\definecolor{matamber}{HTML}{ffc107}
\definecolor{matorange}{HTML}{ff9800}
\definecolor{matdeeporange}{HTML}{ff5722}
\definecolor{matbrown}{HTML}{795548}
\definecolor{matgray}{HTML}{9e9e9e}
\definecolor{matbluegray}{HTML}{607d8b}

% ===========================================
% COLORLET FOR TIKZ
% ===========================================
\colorlet{xcol}{blue!60!black}

% ===========================================
% DARK MODE TCOLORBOX
% ===========================================
\newtcolorbox{qq}[2][]{%
    boxrule=0.75pt,
    sharp corners,
    colframe=white,
    colback=pag,
    coltext=white,
    #1,
}

% Dark definition box
\newtcolorbox{darkdfn}[2][]{%
    enhanced,
    boxrule=1pt,
    sharp corners,
    colframe=p5,
    colback=pag,
    coltext=white,
    coltitle=white,
    fonttitle=\bfseries,
    title=#2,
    #1,
}

% Dark theorem box
\newtcolorbox{darkthm}[2][]{%
    enhanced,
    boxrule=1pt,
    sharp corners,
    colframe=matpurple,
    colback=pag,
    coltext=white,
    coltitle=white,
    fonttitle=\bfseries,
    title=#2,
    #1,
}

% ===========================================
% TIKZ 3D FIX
% ===========================================
% Small fix for canvas is xy plane at z
% https://tex.stackexchange.com/a/48776/121799
\makeatletter
\tikzoption{canvas is xy plane at z}[]{%
    \def\tikz@plane@origin{\pgfpointxyz{0}{0}{#1}}%
    \def\tikz@plane@x{\pgfpointxyz{1}{0}{#1}}%
    \def\tikz@plane@y{\pgfpointxyz{0}{1}{#1}}%
    \tikz@canvas@is@plane}
\makeatother

% ===========================================
% CONDITIONAL LINE TIKZSET
% ===========================================
\tikzset{
    conditional line/.style 2 args={
        /utils/exec={\pgfmathsetmacro{\xcoordA}{#1}},
        /utils/exec={\pgfmathsetmacro{\xcoordB}{#2}},
        insert path={
            \ifdim\xcoordA pt>\xcoordB pt
            (\xcoordA,0) -- (\xcoordB,0)
            \fi
        },
    },
}

% ===========================================
% PGFPLOTS DARK MODE DEFAULTS
% ===========================================
\pgfplotsset{
    dark axis/.style={
        axis line style={white},
        tick style={white},
        label style={white},
        grid style={pag!70},
        every axis label/.append style={white},
        every tick label/.append style={white},
    },
}

% ===========================================
% TIKZ EXTERNALIZATION (for fast compilation)
% ===========================================
% This creates .md5 cache files for figures
% Requires: pdflatex -shell-escape
%
% To enable, add AFTER loading this file:
%   \enabletikzexternalize
%
\newcommand{\enabletikzexternalize}{%
  \usetikzlibrary{external}%
  \tikzexternalize[prefix=figures/]%
}

% ===========================================
% CONVENIENT SHORTCUTS FOR DARK PLOTS
% ===========================================
\newcommand{\darkplot}[1]{%
    \begin{tikzpicture}
    \begin{axis}[
        dark axis,
        grid=both,
        major grid style={line width=.2pt,draw=pag!70},
        #1
    ]
}


% ===========================================
% QUICK GEOMETRY MACROS (tkz-euclide)
% ===========================================
% Draw a point: \point{name}{x}{y}
\newcommand{\gpoint}[3]{\tkzDefPoint(#2,#3){#1}}

% Draw line through points: \gline{A}{B}
\newcommand{\gline}[2]{\tkzDrawLine[color=p4](#1,#2)}

% Draw circle: \gcircle{center}{radius}
\newcommand{\gcircle}[2]{\tkzDrawCircle[color=p5](#1,#2)}

% Mark angle: \gangle{A}{B}{C}
\newcommand{\gangle}[3]{\tkzMarkAngle[color=p3,size=0.5](#1,#2,#3)}

% ===========================================
% USAGE EXAMPLE
% ===========================================
% In your document:
%
% \begin{tikzpicture}
%   \begin{axis}[dark axis, ...]
%     \addplot[p4, thick] {sin(deg(x))};
%   \end{axis}
% \end{tikzpicture}
%
% Or with qq box:
% \begin{qq}{My Title}
%   Content here in dark box
% \end{qq}


\course{Studies-in-Algebra-and-Group-Theory}

\begin{document}

\lecture{16}{Real and Quaternionic Representations}

We have studied actions of finite groups on complex vector spaces. Now we consider the real case and the relationship between real and complex representations.

\section{Real Representations}

Let \(G\) be a finite group. A real representation of \(G\) is a real vector space \(V_0\) together with a group homomorphism \(G \to \GL(V_0)\) (automorphisms as an \(\RR\)-vector space).

The averaging trick works over \(\RR\): starting from any inner product on \(V_0\) and averaging over \(G\), we obtain a \(G\)-invariant (positive-definite, symmetric) inner product. So if \(W_0 \subset V_0\) is a \(G\)-invariant subspace, then \(W_0^\perp\) is also \(G\)-invariant, and every real representation of a finite group is completely reducible.

However, Schur's lemma in its strong form fails over \(\RR\). For instance, \(\ZZ/n\) acts on \(\RR^2\) by rotation by \(2\pi/n\), which is irreducible (for \(n \geq 3\)) but has nontrivial endomorphisms: any rotation commuting with rotation by \(2\pi/n\) works, so \(\End_G(\RR^2) \cong \CC\) (not just \(\RR\)).

\section{Complexification}

The main tool for relating real and complex representations is complexification.

\dfn{Complexification}{Given a real representation \(V_0\) of \(G\), its complexification is the complex vector space \(V = V_0 \otimes_{\RR} {\CC} = V_0 \oplus i V_0\), with \(G\) acting by \(g(v + iw) = gv + i \, gw\) for \(v, w \in V_0\). This defines a functor \(\{\text{real reps of } G\} \to \{\text{complex reps of } G\}\).%
}

\dfn{Real complex representation}{A complex representation \(V\) of \(G\) is called real if there exists a real representation \(V_0\) such that \(V \cong V_0 \otimes_{\RR} {\CC}\).%
}

\mprop{Necessary condition}{If \(V = V_0 \otimes_{\RR} {\CC}\) is real, then \(\chi_V\) takes values in \({\RR}\). (The trace of a matrix with real entries is real.)%
}

The converse is also true for irreducible representations, as we will show.

\section{Characterizing Real Representations}

\mprop{Antilinear involution criterion}{A complex representation \(V\) of \(G\) is real if and only if there exists a \(G\)-equivariant, complex antilinear map \(\tau : V \to V\) (i.e.\ \(\tau(\lambda v) = \ol{\lambda} \tau(v)\) for all \(\lambda \in \CC\)) with \(\tau^2 = \id\).}

\pf{Proof}{(\(\Rightarrow\)) If \(V = V_0 \otimes_{\RR} {\CC}\), define \(\tau(v + iw) = v - iw\) for \(v, w \in V_0\) (complex conjugation with respect to the real form \(V_0\)). Then \(\tau\) is antilinear, \(\tau^2 = \id\), and \(G\)-equivariant since \(g\) preserves \(V_0\).

(\(\Leftarrow\)) Given such a \(\tau\), every \(v \in V\) decomposes as
\[v = \underbrace{\frac{v + \tau(v)}{2}}_{\in V^{\tau}} + \underbrace{\frac{v - \tau(v)}{2}}_{\in V^{-\tau}}\]
where \(V^\tau = \ker(\tau - \id)\) is the \(+1\)-eigenspace and \(V^{-\tau} = \ker(\tau + \id)\) is the \(-1\)-eigenspace. Since \(\tau\) is antilinear, \(\tau(iv) = -i\tau(v)\), so \(V^{-\tau} = iV^\tau\). Thus \(V^\tau\) is an \({\RR}\)-subspace (not a \({\CC}\)-subspace!) with \(V = V^\tau \oplus iV^\tau = V^\tau \otimes_{\RR} {\CC}\). Since \(\tau\) is \(G\)-equivariant, \(V_0 := V^\tau\) is preserved by \(G\), giving a real representation with \(V = V_0 \otimes_{\RR} {\CC}\).%
}


\section{The Frobenius--Schur Indicator}

For irreducible representations, we can determine whether a representation is real, quaternionic, or neither using the bilinear form it carries.

\mprop{Invariant bilinear form from self-duality}{Let \(V\) be an irreducible complex representation of \(G\) with \(\chi_V\) real-valued. Then \(\chi_V = \ol{\chi_V} = \chi_{V^*}\), so \(V \cong V^*\). By Schur's lemma, \(\dim \Hom_G(V, V^*) = 1\), so there is a \(G\)-invariant bilinear form \(B : V \times V \to \CC\), unique up to scaling, and nondegenerate.
}

Since \(B \in (V \otimes V)^* \cong \Sym^2 V^* \oplus \bigwedge^2 V^*\), and both the symmetric and skew-symmetric parts are also \(G\)-invariant, uniqueness forces exactly one of them to be zero. So \(B\) is either symmetric or skew-symmetric. This dichotomy distinguishes real from quaternionic representations.

\mprop{Real iff symmetric form}{An irreducible complex representation \(V\) of a finite group \(G\) is real if and only if \(V\) carries a \(G\)-invariant nondegenerate symmetric bilinear form \(B : V \times V \to \CC\).}

\pf{Proof}{(\(\Rightarrow\)) If \(V = V_0 \otimes_{\RR} {\CC}\), then \(V_0\) carries a \(G\)-invariant real inner product \(B_0\) (by the averaging argument). Extend \({\CC}\)-bilinearly:
\[B(v_1 + iw_1,\, v_2 + iw_2) = B_0(v_1, v_2) + iB_0(w_1, v_2) + iB_0(v_1, w_2) - B_0(w_1, w_2).\]
This is a nondegenerate, symmetric, \(G\)-invariant \({\CC}\)-bilinear form on \(V\).

(\(\Leftarrow\)) Given \(B\), the form determines a \(G\)-equivariant isomorphism \(\varphi : V \xrightarrow{\sim} V^*\). Choose a \(G\)-invariant Hermitian inner product \(H\) on \(V\) (by averaging), which gives a \(G\)-equivariant, antilinear isomorphism \(V \xrightarrow{\sim} V^*\) (via \(v \mapsto H(-, v)\)). Composing one with the inverse of the other yields a \(G\)-equivariant, antilinear map \(\tau : V \to V\), characterized by \(H(\tau(v), w) = B(v, w)\).

Now \(\tau^2 : V \to V\) is \({\CC}\)-linear and \(G\)-equivariant, so by Schur's lemma, \(\tau^2 = \lambda \cdot \id\) for some \(\lambda \in {\CC}\). We show \(\lambda \in {\RR}_{> 0}\):
\[H(\tau^2(v), v) = B(\tau(v), v) = B(v, \tau(v)) = H(\tau(v), \tau(v)) \geq 0\]
where the second equality uses symmetry of \(B\). So \(\lambda \geq 0\), and \(\lambda > 0\) since \(\tau \neq 0\). Replacing \(\tau\) by \(\lambda^{-1/2} \tau\) (which is still antilinear and \(G\)-equivariant), we arrange \(\tau^2 = \id\). By the antilinear involution criterion, \(V\) is real.
}

\section{Quaternionic Representations}

In the other case, the invariant bilinear form \(B\) is skew-symmetric. The same argument produces a \(G\)-equivariant antilinear map \(J : V \to V\) with \(J^2 = -\id\).

\dfn{Quaternionic structure}{A quaternionic structure on a complex vector space \(V\) is a \(\CC\)-antilinear map \(J : V \to V\) with \(J^2 = -\id\). A representation with such a \(G\)-equivariant \(J\) is called a quaternionic representation.}

Such a \(J\) makes \(V\) into a module over the quaternions \(\mathbb{H} = \{a + bi + cj + dk \mid a, b, c, d \in \RR\}\), the division algebra (noncommutative ring where every nonzero element has a multiplicative inverse) with \(i^2 = j^2 = k^2 = ijk = -1\). Concretely, \(\mathbb{H} = \CC \oplus \CC j\) with \(ji = -ij\), \(j^2 = -1\), so an \(\mathbb{H}\)-module is a \(\CC\)-vector space with an antilinear map \(J\) (``right multiplication by \(j\)'') satisfying \(J^2 = -\id\).

\nt{A quaternionic representation has \(\chi_V\) real-valued (since \(V \cong V^*\)) but does \emph{not} come from a real representation. The three cases for an irreducible \(V\) with real character are:
\begin{itemize}
    \item \(V\) admits a \(G\)-invariant nondegenerate symmetric bilinear form \(\iff\) \(V\) is real \(\iff\) \(V = V_0 \otimes_{\RR} {\CC}\).
    \item \(V\) admits a \(G\)-invariant nondegenerate skew-symmetric bilinear form \(\iff\) \(V\) is quaternionic \(\iff\) \(V\) carries a \(G\)-equivariant antilinear \(J\) with \(J^2 = -\id\).
\end{itemize}
And if \(\chi_V\) is not real-valued, then \(V \not\cong V^*\) and \(V\) is neither real nor quaternionic.}

\section{The Frobenius--Schur Indicator}

The symmetric/skew dichotomy can be detected numerically.

\thm{Frobenius--Schur indicator}{For an irreducible representation \(V\) of a finite group \(G\),
\[\nu(V) := \frac{1}{|G|} \sum_{g \in G} \chi_V(g^2) = \begin{cases} 1 & \text{if } V \text{ is real}, \\ 0 & \text{if } \chi_V \notin {\RR}, \\ -1 & \text{if } V \text{ is quaternionic}. \end{cases}\]%
}

\pf{Proof}{Using the character formulas \(\chi_{\Sym^2 V}(g) = \frac{1}{2}(\chi_V(g)^2 + \chi_V(g^2))\) and \(\chi_{\bigwedge^2 V}(g) = \frac{1}{2}(\chi_V(g)^2 - \chi_V(g^2))\), we get
\[\nu(V) = \frac{1}{|G|} \sum_g \chi_V(g^2) = \dim(\Sym^2 V)^G - \dim(\bigwedge^2 V)^G.\]
Now \((\Sym^2 V)^G\) counts the multiplicity of the trivial representation in \(\Sym^2 V\), and similarly for \(\bigwedge^2 V\). A \(G\)-invariant bilinear form on \(V\) (a \(G\)-fixed element of \((V \otimes V)^*\)) is either symmetric or skew-symmetric.

If \(V \not\cong V^*\), then \(\Hom_G(V, V^*) = 0\), so no invariant bilinear form exists. Both dimensions are 0, giving \(\nu(V) = 0\).

If \(V \cong V^*\) (i.e., \(\chi_V\) real-valued), then \(\dim \Hom_G(V, V^*) = 1\) by Schur. The unique form is either symmetric or skew, so exactly one of the two dimensions equals 1. The form is symmetric iff \(V\) is real (\(\nu = 1 - 0 = 1\)), and skew iff \(V\) is quaternionic (\(\nu = 0 - 1 = -1\)).%
}

\mprop{Quaternionic representations have even dimension}{If \(V\) is quaternionic, then \(\dim_{\CC} V\) is even, since \(J\) with \(J^2 = -\id\) makes \(V\) into an \(\mathbb{H}\)-module and \(\dim_{\CC} V = 2 \dim_\mathbb{H} V\). In particular, an irreducible representation with real character and odd dimension is automatically real.%
}


\section{Examples}

\ex{Standard representation of \(S_3\)}{The 2-dimensional irreducible \(V\) of \(S_3 \cong D_3\) is real: \(S_3\) acts on \(V_0 = \{(x_1, x_2, x_3) \in {\RR}^3 : \sum x_i = 0\} \cong {\RR}^2\) by rotations and reflections, and \(V = V_0 \otimes_{\RR} {\CC}\).

Alternatively: \(V^* \cong V\) (since \(\chi_V\) is real-valued), and \(\bigwedge^2 V^* \cong U'\) contains no trivial summand, so the \(G\)-invariant bilinear form must be in \(\Sym^2 V^*\). Since \(\Sym^2 V^* \cong U \oplus V\) contains a trivial summand, an invariant symmetric form exists, confirming \(V\) is real.%
}

\ex{The quaternion group \(Q_8\)}{The quaternion group \(Q_8 = \{\pm 1, \pm i, \pm j, \pm k\}\) with \(i^2 = j^2 = k^2 = ijk = -1\) acts faithfully on \({\CC}^2\) by
\[\pm 1 \mapsto \pm I, \quad \pm i \mapsto \pm\begin{pmatrix} i & 0 \\ 0 & -i \end{pmatrix}, \quad \pm j \mapsto \pm\begin{pmatrix} 0 & 1 \\ -1 & 0 \end{pmatrix}, \quad \pm k \mapsto \pm\begin{pmatrix} 0 & i \\ i & 0 \end{pmatrix}.\]
The character takes values \(\chi(\pm 1) = \pm 2\) and \(\chi = 0\) on all other elements---all real. However, this representation is \emph{not} real: it is quaternionic.

To see this cannot come from a real representation: if \(Q_8\) acted faithfully on \({\RR}^2\), the invariant inner product would give \(Q_8 \hookrightarrow O(2)\). But \(O(2)\) has only 2 elements squaring to \(-I\) (rotations by \(\pm 90^\circ\)), while \(Q_8\) needs 6 elements (\(\pm i, \pm j, \pm k\)) with square \(-1\). Contradiction.

The quaternionic structure is: view \({\CC}^2 = \mathbb{H}\) via \((z_1, z_2) \leftrightarrow z_1 + jz_2\), and the above matrices correspond to left multiplication by \(\pm 1, \pm i, \pm j, \pm k\). The antilinear map \(J : V \to V\) with \(J^2 = -\id\) is right multiplication by \(j\), which commutes with left multiplication.%
}


\mer{Exercises}{
\begin{enumerate}[(1)]
    \item Classify each irreducible representation of \(S_4\) as real, quaternionic, or complex-type. (Hint: all characters of \(S_4\) are real-valued, so each is either real or quaternionic. Use the dimension criterion: a quaternionic representation must have even dimension.)

    \item Show that if \(V\) is a quaternionic representation of \(G\) (carrying an antilinear \(J\) with \(J^2 = -\id\)), then \(\dim_{\CC} V\) is even.

    \item Prove the Frobenius--Schur indicator formula: for an irreducible representation \(V\),
    \[\frac{1}{|G|} \sum_{g \in G} \chi_V(g^2) = \begin{cases} 1 & \text{if } V \text{ is real}, \\ 0 & \text{if } \chi_V \notin {\RR}, \\ -1 & \text{if } V \text{ is quaternionic}. \end{cases}\]
    (Hint: \(\frac{1}{|G|}\sum_g \chi_V(g^2) = \dim(\Sym^2 V)^G - \dim(\bigwedge^2 V)^G\), and exactly one of these dimensions is 1 when \(V \cong V^*\).)

    \item Using the Frobenius--Schur indicator, verify that the 2-dimensional irreducible of \(Q_8\) is quaternionic.

    \item Show that the representations \(Y\) and \(Z\) of \(A_5\) (the 3-dimensional irreducibles with golden-ratio character values) are real representations. (Their characters are real, and they have odd dimension.)
\end{enumerate}
}

\begin{solution}
\textbf{(1)} All characters of \(S_4\) are real-valued (check the character table). A quaternionic representation has \(J : V \to V\) with \(J^2 = -\id\), and \(J\) pairs basis vectors into orbits of size 2, so \(\dim V\) must be even. The irreducibles of \(S_4\) have dimensions 1, 1, 2, 3, 3. The odd-dimensional ones (\(U, U', V_4, V_4'\)) are automatically real. For \(W\) (\(\dim = 2\)): compute the Frobenius--Schur indicator \(\frac{1}{24}\sum \chi_W(g^2)\). We have \(\chi_W = (2, 0, 2, -1, 0)\) with class sizes \((1, 6, 3, 8, 6)\). Squaring: \(e \to e, (12) \to e, (12)(34) \to e, (123) \to (132), (1234) \to (13)(24)\). So \(\chi_W(g^2) = (2, 2, 2, -1, 2)\). The indicator is \(\frac{1}{24}(2 + 12 + 6 - 8 + 12) = 1\). So \(W\) is also real.

\textbf{(2)} The antilinear map \(J\) gives an \({\RR}\)-linear map \(V \to V\) whose square is \(-\id\). Viewing \(V\) as \({\RR}^{2n}\) (where \(\dim_{\CC} V = n\)), this gives a real matrix \(M\) with \(M^2 = -I\). The characteristic polynomial \(\det(M - tI)\) has degree \(2n\) and all roots are \(\pm i\), which come in conjugate pairs. More directly: \(J\) makes \(V\) into an \(\mathbb{H}\)-module, and \(\dim_\mathbb{H} V = \frac{1}{2}\dim_{\CC} V\) must be a positive integer.

\textbf{(3)} We have \(\frac{1}{|G|}\sum_g \chi_V(g^2) = \frac{1}{|G|}\sum_g \Tr(\rho(g^2))\). Using the decomposition \(V \otimes V = \Sym^2 V \oplus \bigwedge^2 V\) and the identities \(\chi_{\Sym^2 V}(g) = \frac{1}{2}(\chi_V(g)^2 + \chi_V(g^2))\), \(\chi_{\bigwedge^2 V}(g) = \frac{1}{2}(\chi_V(g)^2 - \chi_V(g^2))\), we get the indicator \(= \dim(\Sym^2 V)^G - \dim(\bigwedge^2 V)^G\). When \(V \cong V^*\), the unique invariant bilinear form is either symmetric (contributing 1 to \(\Sym^2\)) or skew (contributing 1 to \(\bigwedge^2\)). When \(V \not\cong V^*\), neither \(\Sym^2 V\) nor \(\bigwedge^2 V\) contains a trivial summand, giving 0.

\textbf{(4)} \(Q_8\) has 5 conjugacy classes: \(\{1\}, \{-1\}, \{\pm i\}, \{\pm j\}, \{\pm k\}\) of sizes 1, 1, 2, 2, 2. For the 2-dimensional rep: \(\chi = (2, -2, 0, 0, 0)\). Squaring: \(1^2 = 1\), \((-1)^2 = 1\), \((\pm i)^2 = -1\), \((\pm j)^2 = -1\), \((\pm k)^2 = -1\). So \(\chi(g^2) = (2, 2, -2, -2, -2)\). Indicator: \(\frac{1}{8}(2 + 2 - 4 - 4 - 4) = -1\). Quaternionic.

\textbf{(5)} Both \(Y\) and \(Z\) have real characters and odd dimension 3. Since a quaternionic representation must have even complex dimension (exercise 2), odd-dimensional representations with real characters are automatically real.
\end{solution}

\end{document}
